\chapter{2018年浙江卷}

\section{单选题}

\begin{problem}
已知全集 \(U=\{1,2,3,4,5\}\)，\(A=\{1,3\}\)，则 \(\complement_U A=\)（\quad）

\choices
    {\(\varnothing\)}
    {\(\{1,3\}\)}
    {\(\{2,4,5\}\)}
    {\(\{1,2,3,4,5\}\)}

\end{problem}

\begin{answer}
C
\end{answer}

\begin{solution}
由补集的定义可知，\(\complement_U A\) 是由所有属于 \(U\) 但不属于 \(A\) 的元素组成的集合，所以 \(\complement_U A=\{2,4,5\}\).
故选 C.
\end{solution}
\begin{problem}
双曲线 \(\dfrac{x^2}{3}-y^2=1\) 的焦点坐标是（\quad）

\choices
    {\((-\sqrt2,0)\)，\((\sqrt2,0)\)}
    {\((-2,0)\)，\((2,0)\)}
    {\((0,-\sqrt2)\)，\((0,\sqrt2)\)}
    {\((0,-2)\)，\((0,2)\)}

\end{problem}

\begin{answer}
B
\end{answer}

\begin{solution}
由双曲线的标准方程可知，其焦点在 \(x\) 轴上，且 \(a^2=3\)，\(b^2=1\)，所以 \(c=\sqrt{a^2+b^2}=\sqrt4=2\).
因此焦点坐标为 \((-2,0)\) 和 \((2,0)\). 故选 B.
\end{solution}
\begin{problem}
某几何体的三视图如图所示（单位：cm），则该几何体的体积（单位：cm\(^3\)）是（\quad）

\bitmapfigure[max width=0.58\linewidth,max totalheight=4.8cm]{img/2018/zhejiang/q3_fig1.png}

\choices
    {\(2\)}
    {\(4\)}
    {\(6\)}
    {\(8\)}

\end{problem}

\begin{answer}
C
\end{answer}

\begin{solution}
由三视图可知，该几何体是一个底面为直角梯形的四棱柱. 由图可得其底面面积为 \(3\)，高为 \(2\)，所以体积为 \(V=3\times2=6\).
故选 C.
\end{solution}
\begin{problem}
复数 \(\dfrac{2}{1-\i}\)（\(\i\) 为虚数单位）的共轭复数是（\quad）

\choices
    {\(1+\i\)}
    {\(1-\i\)}
    {\(-1+\i\)}
    {\(-1-\i\)}

\end{problem}

\begin{answer}
B
\end{answer}

\begin{solution}
先化简 \(\frac{2}{1-\i}=\frac{2(1+\i)}{(1-\i)(1+\i)}=1+\i\).
所以它的共轭复数是 \(1-\i\). 故选 B.
\end{solution}
\begin{problem}
函数 \(y=2^{|x|}\sin 2x\) 的图象可能是（\quad）

\choices
    {\(\bitmapinclude[max width=0.17\textwidth,max totalheight=4.8cm]{img/2018/zhejiang/q5_choice_a.png}\)}
    {\(\bitmapinclude[max width=0.17\textwidth,max totalheight=4.8cm]{img/2018/zhejiang/q5_choice_b.png}\)}
    {\(\bitmapinclude[max width=0.17\textwidth,max totalheight=4.8cm]{img/2018/zhejiang/q5_choice_c.png}\)}
    {\(\bitmapinclude[max width=0.17\textwidth,max totalheight=4.8cm]{img/2018/zhejiang/q5_choice_d.png}\)}

\end{problem}

\begin{answer}
D
\end{answer}

\begin{solution}
函数 \(y=2^{|x|}\sin 2x\) 是奇函数，图象关于原点对称，故排除 A、B. 又当 \(x\) 取正负较大值时，\(|y|\) 的振幅增大，图象应与 D 一致. 故选 D.
\end{solution}
\begin{problem}
已知平面 \(\alpha\)，直线 \(m,n\) 满足 \(m\not\subset \alpha\)，\(n\subset \alpha\)，则“\(m\parallel n\)”是“\(m\parallel \alpha\)”的（\quad）

\choices
    {充分不必要条件}
    {必要不充分条件}
    {充分必要条件}
    {既不充分也不必要条件}

\end{problem}

\begin{answer}
A
\end{answer}

\begin{solution}
若 \(m\parallel n\) 且 \(n\subset\alpha\)，则 \(m\) 与平面 \(\alpha\) 内一条直线平行，所以 \(m\parallel \alpha\)，充分性成立. 反过来，\(m\parallel \alpha\) 时，不一定有 \(m\parallel n\)，所以必要性不成立. 故选 A.
\end{solution}
\begin{problem}
设 \(0<p<1\)，随机变量 \(\xi\) 的分布列是

\begin{center}
\begin{tabular}{c|ccc}
\(\xi\) & \(0\) & \(1\) & \(2\) \\
\hline
\(P\) & \(\dfrac{1-p}{2}\) & \(\dfrac12\) & \(\dfrac p2\)
\end{tabular}
\end{center}

则当 \(p\) 在 \((0,1)\) 内增大时，\(D(\xi)\)（\quad）

\choices
    {减小}
    {增大}
    {先减小后增大}
    {先增大后减小}

\end{problem}

\begin{answer}
D
\end{answer}

\begin{solution}
由分布列得 \(E(\xi)=\frac12+p\)，\(E(\xi^2)=\frac12+2p\). 所以 \(D(\xi)=E(\xi^2)-[E(\xi)]^2=\frac14+p-p^2\).
其导数为 \(1-2p\)，故在 \((0,1)\) 内先增大后减小. 故选 D.
\end{solution}
\begin{problem}
已知四棱锥 \(S-ABCD\) 的底面是正方形，侧棱长均相等，\(E\) 是线段 \(AB\) 上的点（不含端点），设 \(SE\) 与 \(BC\) 所成的角为 \(\theta_1\)，\(SE\) 与平面 \(ABCD\) 所成的角为 \(\theta_2\)，二面角 \(S-AB-C\) 的平面角为 \(\theta_3\)，则（\quad）

\bitmapfigure[max width=0.52\linewidth,max totalheight=4.8cm]{img/2018/zhejiang/q8_fig1.png}

\choices
    {\(\theta_1\le \theta_2\le \theta_3\)}
    {\(\theta_3\le \theta_2\le \theta_1\)}
    {\(\theta_1\le \theta_3\le \theta_2\)}
    {\(\theta_2\le \theta_3\le \theta_1\)}

\end{problem}

\begin{answer}
D
\end{answer}

\begin{solution}
由正方形底面和侧棱相等可知，\(S\) 在底面上的射影是正方形中心. 比较直线 \(SE\) 在底面内的射影、与底面的夹角以及二面角的平面角，可得 \(\theta_2\le\theta_3\le\theta_1\).
故选 D.
\end{solution}
\begin{problem}
已知 \(a,b,e\) 是平面向量，\(e\) 是单位向量. 若非零向量 \(a\) 与 \(e\) 的夹角为 \(\dfrac{\pi}{3}\)，向量 \(b\) 满足 \(\lVert b\rVert^2-4e\cdot b+3=0\)，则 \(\lVert a-b\rVert\) 的最小值是（\quad）

\choices
    {\(\sqrt3-1\)}
    {\(\sqrt3+1\)}
    {\(2\)}
    {\(2-\sqrt3\)}

\end{problem}

\begin{answer}
A
\end{answer}

\begin{solution}
由 \(\lVert b\rVert^2-4e\cdot b+3=0\) 得 \(\lVert b-2e\rVert^2=1\)，
所以 \(b\) 的端点在以 \(2e\) 为圆心、1 为半径的圆上. 又 \(a\) 与 \(e\) 的夹角为 \(\dfrac{\pi}{3}\)，结合图形可知，\(\lVert a-b\rVert\) 的最小值为 \(\sqrt3-1\). 故选 A.
\end{solution}
\begin{problem}
已知 \(a_1,a_2,a_3,a_4\) 成等比数列，且 \(a_1+a_2+a_3+a_4=\ln(a_1+a_2+a_3)\). 若 \(a_1>1\)，则（\quad）

\choices
    {\(a_1<a_3\)，\(a_2<a_4\)}
    {\(a_1>a_3\)，\(a_2<a_4\)}
    {\(a_1<a_3\)，\(a_2>a_4\)}
    {\(a_1>a_3\)，\(a_2>a_4\)}

\end{problem}

\begin{answer}
B
\end{answer}

\begin{solution}
设等比公比为 \(q\). 由于 \(a_1>1\) 且 \(a_1+a_2+a_3+a_4=\ln(a_1+a_2+a_3)\)，可知 \(q\in(-1,0)\). 因此奇数项同号、偶数项同号，且 \(a_1>a_3\)，\(a_2<a_4\). 故选 B.
\end{solution}
\section{填空题}

\begin{problem}
已知数列 \(\{a_n\}\) 满足 \(a_n=\dfrac{n(n+1)}{2}\)，则 \(S_3=\fillinblank{}\).

\end{problem}

\begin{answer}
\(10\)
\end{answer}

\begin{solution}
因为 \(a_1=1\)，\(a_2=3\)，\(a_3=6\)，所以 \(S_3=1+3+6=10\).
\end{solution}
\begin{problem}
设圆 \(C\) 的圆心坐标是 \((0,m)\)，半径长是 \(r\). 若直线 \(2x-y+3=0\) 与圆相切于点 \(A(-2,-1)\)，则 \(m=\fillinblank{}\)，\(r=\fillinblank{}\).

\end{problem}

\begin{answer}
\(-2\)，\(\sqrt5\)
\end{answer}

\begin{solution}
由切点 \(A(-2,-1)\) 可知，半径 \(OA\) 垂直于直线 \(2x-y+3=0\)，其斜率为 \(-\dfrac12\). 过 \(A\) 作垂线得圆心在直线 \(y+1=-\frac12(x+2)\) 上，代入 \(x=0\) 得 \(m=-2\).

半径为 \(r=OA=\sqrt{(-2)^2+(-1+2)^2}=\sqrt5\).
\end{solution}
\begin{problem}
设 \((\sqrt2+x)^9=a_1+a_2x+a_3x^2+\cdots+a_{10}x^9\)，则常数项是 \fillinblank{}；系数为有理数的项的个数是 \fillinblank{}.

\end{problem}

\begin{answer}
\(16\sqrt2\)，\(5\)
\end{answer}

\begin{solution}
展开式的通项为 \(T_{r+1}=C_9^r(\sqrt2)^{9-r}x^r\)（\(r=0,1,\dots,9\)）. 常数项对应 \(r=0\)，故常数项为 \((\sqrt2)^9=16\sqrt2\).
系数为有理数时，\(9-r\) 为偶数，所以 \(r\) 为奇数，共有 \(5\) 项.
\end{solution}
\begin{problem}
已知 \(\tan\theta=2\)，则 \(\cos2\theta=\fillinblank{}\)；\(\tan\left(\theta-\dfrac{\pi}{4}\right)=\fillinblank{}\).

\end{problem}

\begin{answer}
\(-\dfrac35\)，\(\dfrac13\)
\end{answer}

\begin{solution}
由 \(\cos2\theta=\frac{1-\tan^2\theta}{1+\tan^2\theta}=\frac{1-4}{1+4}=-\frac35\)，又 \(\tan\left(\theta-\frac{\pi}{4}\right)=\frac{\tan\theta-1}{1+\tan\theta}=\frac{2-1}{1+2}=\frac13\).
\end{solution}
\begin{problem}
二项式 \(\left(\sqrt{x}+\dfrac{1}{2x}\right)^8\) 的展开式的常数项是 \fillinblank{}.

\end{problem}

\begin{answer}
\(0\)
\end{answer}

\begin{solution}
展开式的通项为 \(T_{r+1}=C_8^r(\sqrt{x})^{8-r}\left(\frac{1}{2x}\right)^r=C_8^r2^{-r}x^{4-\frac{3r}{2}}\). 令指数 \(4-\dfrac{3r}{2}=0\)，得 \(r=\dfrac83\)，不是整数，所以展开式中没有常数项，常数项为 \(0\).
\end{solution}
\begin{problem}
已知 \(\lambda\in\R\)，函数 \(f(x)=x-4\)（\(x\ge\lambda\)），\(f(x)=x^2-4x+3\)（\(x<\lambda\)）. 当 \(\lambda=2\) 时，不等式 \(f(x)<0\) 的解集是 \fillinblank{}. 若函数 \(f(x)\) 恰有 2 个零点，则 \(\lambda\) 的取值范围是 \fillinblank{}.

\end{problem}

\begin{answer}
\((1,4)\)，\((1,3]\cup(4,+\infty)\)
\end{answer}

\begin{solution}
当 \(\lambda=2\) 时，
\examdisplaycases{f(x)=}{
x-4, & x\ge2,\\
x^2-4x+3, & x<2.
}
于是 \(x\ge2\) 时，\(f(x)<0\) 等价于 \(x<4\)，得到 \([2,4)\)；\(x<2\) 时，\(f(x)<0\) 等价于 \(1<x<2\). 故解集为 \((1,4)\).

再看零点. 二次部分的零点为 \(x=1,3\)，一次部分的零点为 \(x=4\). 当且仅当 \(1<\lambda\le3\) 或 \(\lambda>4\) 时，恰有两个零点.
\end{solution}
\begin{problem}
从 \(1,3,5,7,9\) 中任取 2 个数字，从 \(0,2,4,6\) 中任取 2 个数字，一共可以组成 \fillinblank{} 个没有重复数字的四位数. （用数字作答）

\end{problem}

\begin{answer}
\(1260\)
\end{answer}

\begin{solution}
分两类讨论.

不含 \(0\) 时，从 \(1,3,5,7,9\) 中选 2 个有 \(C_5^2=10\) 种，从 \(2,4,6\) 中选 2 个有 \(C_3^2=3\) 种，每组可排成 \(4!=24\) 个四位数，共有 \(10\times3\times24=720\) 个.

含 \(0\) 时，从 \(1,3,5,7,9\) 中选 2 个有 \(10\) 种，从 \(2,4,6\) 中选 1 个有 \(3\) 种，再把这 4 个数字排列，其中千位不能是 \(0\)，共有 \(3\times 18=540\) 个.

总数为
\(720+540=1260\).
\end{solution}
\begin{problem}
已知点 \(P(0,1)\)，椭圆 \(\dfrac{x^2}{4}+y^2=m\)（\(m>1\)） 上两点 \(A,B\) 满足 \(\overrightarrow{PA}=2\overrightarrow{PB}\)，则当 \(m=\fillinblank{}\) 时，点 \(B\) 横坐标的绝对值最大.

\end{problem}

\begin{answer}
\(5\)
\end{answer}

\begin{solution}
设 \(A(x_1,y_1)\)，\(B(x_2,y_2)\).

由 \(P(0,1)\)，\(\overrightarrow{PA}=2\overrightarrow{PB}\)，

可得 \(-x_1=2x_2\)，\(1-y_1=2(y_2-1)\)，

即有 \(x_1=-2x_2\)，\(y_1+2y_2=3\).

又 \(x_1^2+4y_1^2=4m\)，

即为 \(x_2^2+y_1^2=m\)，\circled{1}

即有 \(x_2^2+4y_2^2=4m\). \circled{2}

\circled{1} \(-\) \circled{2} 得 \((y_1-2y_2)(y_1+2y_2)=-3m\)，

可得 \(y_1-2y_2=-m\).

解得 \(y_1=\frac{3-m}{2}\)，\(y_2=\frac{m+3}{4}\).

则 \(m=x_2^2+\left( \frac{3-m}{2} \right)^2\)，

即有 \(x_2^2=m-\left( \frac{3-m}{2} \right)^2=\frac{-m^2+10m-9}{4}\).

即有 \(m=5\) 时，\(x_2^2\) 有最大值 \(4\)，

即点 \(B\) 横坐标的绝对值最大.

故答案为：\(5\).
\end{solution}
\section{解答题}

\begin{problem}
已知角 \(\alpha\) 的顶点与原点 \(O\) 重合，始边与 \(x\) 轴的非负半轴重合，它的终边过点 \(P\left(-\dfrac35,-\dfrac45\right)\).

\begin{enumerate}
  \item 求 \(\sin(\alpha+\pi)\) 的值；
  \item 若角 \(\beta\) 满足 \(\sin(\alpha+\beta)=\dfrac{5}{13}\)，求 \(\cos\beta\) 的值.
\end{enumerate}

\end{problem}

\begin{solution}
（1）由角 \(\alpha\) 的终边过点 \(P\left( -\dfrac{3}{5},-\dfrac{4}{5} \right)\) 得
\(\sin \alpha=-\frac{4}{5}\)，所以 \(\sin(\alpha+\pi)=-\sin \alpha=\frac{4}{5}\).

（2）由角 \(\alpha\) 的终边过点 \(P\left( -\dfrac{3}{5},-\dfrac{4}{5} \right)\) 得
由 \(\sin(\alpha+\beta)=\dfrac{5}{13}\) 得 \(\cos \alpha=-\frac{3}{5}\)，\(\cos(\alpha+\beta)=\pm\frac{12}{13}\).

由 \(\beta=(\alpha+\beta)-\alpha\) 得
\(\cos \beta=\cos(\alpha+\beta)\cos \alpha+\sin(\alpha+\beta)\sin \alpha\)，所以 \(\cos \beta=-\frac{56}{65}\)或\(\cos \beta=-\frac{16}{65}\).
\end{solution}
\begin{problem}
如图，已知多面体 \(ABCA_1B_1C_1\)，\(A_1A\)，\(B_1B\)，\(C_1C\) 均垂直于平面 \(ABC\)，\(\angle ABC=120^\circ\)，\(A_1A=4\)，\(C_1C=1\)，\(AB=BC=B_1B=2\).

\begin{enumerate}
  \item 证明：\(AB_1\perp\) 平面 \(A_1B_1C_1\)；
  \item 求直线 \(AC_1\) 与平面 \(ABB_1\) 所成的角的正弦值.
\end{enumerate}

\bitmapfigure[max width=0.42\linewidth,max totalheight=4.8cm]{img/2018/zhejiang/q19_fig1.png}

\end{problem}

\begin{solution}
方法一：

（1）由 \(AB=2\)，\(AA_1=4\)，\(BB_1=2\)，\(AA_1 \perp AB\)，\(BB_1 \perp AB\) 得 \(AB_1=A_1B_1=2\sqrt{2}\)，
所以
\(A_1B_1^2+AB_1^2=AA_1^2\).
故 \(AB_1 \perp A_1B_1\).

由 \(BC=2\)，\(BB_1=2\)，\(CC_1=1\)，\(BB_1 \perp BC\)，\(CC_1 \perp BC\) 得 \(B_1C_1=\sqrt{5}\)，由 \(AB=BC=2\)，\(\angle ABC=120^\circ\) 得 \(AC=2\sqrt{3}\)，由 \(CC_1 \perp AC\)，得 \(AC_1=\sqrt{13}\)，
所以
\(AB_1^2+B_1C_1^2=AC_1^2\)，
故 \(AB_1 \perp B_1C_1\).

因此 \(AB_1 \perp\) 平面 \(A_1B_1C_1\).

（2）如图，过点 \(C_1\) 作 \(C_1D \perp A_1B_1\)，交直线 \(A_1B_1\) 于点 \(D\)，连结 \(AD\).

由 \(AB_1 \perp\) 平面 \(A_1B_1C_1\) 得平面 \(A_1B_1C_1 \perp\) 平面 \(ABB_1\).

由 \(C_1D \perp A_1B_1\) 得 \(C_1D \perp\) 平面 \(ABB_1\).

所以 \(\angle C_1AD\) 是 \(AC_1\) 与平面 \(ABB_1\) 所成的角.

由 \(B_1C_1=\sqrt{5}\)，\(A_1B_1=2\sqrt{2}\)，\(A_1C_1=\sqrt{21}\) 得 \(\cos \angle C_1A_1B_1=\frac{\sqrt{6}}{\sqrt{7}}\)，\(\sin \angle C_1A_1B_1=\frac{1}{\sqrt{7}}\)，所以 \(C_1D=\sqrt{3}\)，
故
\(\sin \angle C_1AD=\frac{C_1D}{AC_1}=\frac{\sqrt{39}}{13}\).

因此，直线 \(AC_1\) 与平面 \(ABB_1\) 所成的角的正弦值是 \(\dfrac{\sqrt{39}}{13}\).

方法二：

（1）如图，以 \(AC\) 的中点 \(O\) 为原点，分别以射线 \(OB\)，\(OC\) 为 \(x\)，\(y\) 轴的正半轴，建立空间直角坐标系 \(O\)-\(xyz\).

由题意知各点坐标如下：
\(A(0,-\sqrt{3},0)\)，\(B(1,0,0)\)，\(A_1(0,-\sqrt{3},4)\)，\(B_1(1,0,2)\)，\(C_1(0,\sqrt{3},1)\)，
因此
 \(\overrightarrow{AB_1}=(1,\sqrt{3},2)\)，\(\overrightarrow{A_1B_1}=(1,\sqrt{3},-2)\)，\(\overrightarrow{A_1C_1}=(0,2\sqrt{3},-3)\)，
由
 \(\overrightarrow{AB_1}\cdot \overrightarrow{A_1B_1}=0\)
得 \(AB_1 \perp A_1B_1\).

由
 \(\overrightarrow{AB_1}\cdot \overrightarrow{A_1C_1}=0\)
得 \(AB_1 \perp A_1C_1\).

所以 \(AB_1 \perp\) 平面 \(A_1B_1C_1\).

（2）设直线 \(AC_1\) 与平面 \(ABB_1\) 所成的角为 \(\theta\).

由（1）可知 \(\overrightarrow{AC_1}=(0,2\sqrt{3},1)\)，\(\overrightarrow{AB}=(1,\sqrt{3},0)\)，\(\overrightarrow{BB_1}=(0,0,2)\)，
设平面 \(ABB_1\) 的法向量 \(\mathbf n=(x,y,z)\).

由
 \(\begin{cases}
\mathbf n\cdot \overrightarrow{AB}=0,\\
\mathbf n\cdot \overrightarrow{BB_1}=0,
\end{cases}\)
即
 \(\begin{cases}
x+\sqrt{3}y=0,\\
2z=0,
\end{cases}\)
可取 \(\mathbf n=(-\sqrt{3},1,0)\).

所以
\(\sin \theta=\left| \cos \angle\left( \overrightarrow{AC_1},\mathbf n \right) \right|=\frac{\left| \overrightarrow{AC_1}\cdot \mathbf n \right|}{\left| \overrightarrow{AC_1} \right|\cdot \left| \mathbf n \right|}=\frac{\sqrt{39}}{13}\).

因此，直线 \(AC_1\) 与平面 \(ABB_1\) 所成的角的正弦值是 \(\dfrac{\sqrt{39}}{13}\).
\end{solution}
\begin{problem}
已知等比数列 \(\{a_n\}\) 的公比 \(q>1\)，且 \(a_3+a_4+a_5=28\)，\(a_4+2\) 是 \(a_3,a_5\) 的等差中项. 数列 \(\{b_n\}\) 满足 \(b_1=1\)，数列 \(\{(b_{n+1}-b_n)a_n\}\) 的前 \(n\) 项和为 \(2n^2+n\).

\begin{enumerate}
  \item 求 \(q\) 的值；
  \item 求数列 \(\{b_n\}\) 的通项公式.
\end{enumerate}

\end{problem}

\begin{solution}
（1）由 \(a_4+2\) 是 \(a_3,a_5\) 的等差中项得
 \(a_3+a_5=2a_4+4\)，
所以
 \(a_3+a_4+a_5=3a_4+4=28\)，
解得
\(a_4=8\).

由 \(a_3+a_5=20\) 得
\(8\left( q+\frac{1}{q} \right)=20\)，因为 \(q>1\)，所以 \(q=2\).

（2）设
 \(c_n=(b_{n+1}-b_n)a_n\)，
数列 \(\{c_n\}\) 前 \(n\) 项和为 \(S_n\).

由
\examdisplaycases{c_n=}{
S_1, & n=1,\\
S_n-S_{n-1}, & n\ge2,
}
解得
\(c_n=4n-1\).

由（1）可知
 \(a_n=2^{n-1}\)，
所以
 \(b_{n+1}-b_n=(4n-1)\left( \frac{1}{2} \right)^{n-1}\)，
故
 \(b_n-b_{n-1}=(4n-5)\left( \frac{1}{2} \right)^{n-2}\)，\(n\ge2\).

\examdisplayaligned{
b_n-b_1&=(b_n-b_{n-1})+(b_{n-1}-b_{n-2})+\cdots+(b_3-b_2)+(b_2-b_1)\\
&=(4n-5)\left( \frac{1}{2} \right)^{n-2}+(4n-9)\left( \frac{1}{2} \right)^{n-3}+\cdots+7\cdot \frac{1}{2}+3.
}

设 \(T_n=3+7\cdot \frac{1}{2}+11\cdot \left( \frac{1}{2} \right)^2+\cdots+(4n-5)\cdot \left( \frac{1}{2} \right)^{n-2}\)，\(n\ge2\). \(\frac{1}{2}T_n=3\cdot \frac{1}{2}+7\cdot \left( \frac{1}{2} \right)^2+\cdots+(4n-9)\cdot \left( \frac{1}{2} \right)^{n-2}+(4n-5)\cdot \left( \frac{1}{2} \right)^{n-1}\)，所以 \(\frac{1}{2}T_n=3+4\cdot \frac{1}{2}+4\cdot \left( \frac{1}{2} \right)^2+\cdots+4\cdot \left( \frac{1}{2} \right)^{n-2}-(4n-5)\cdot \left( \frac{1}{2} \right)^{n-1}\)，因此
\(T_n=14-(4n+3)\cdot \left( \frac{1}{2} \right)^{n-2}\)，\(n\ge2\).

又 \(b_1=1\)，所以
\(b_n=15-(4n+3)\cdot \left( \frac{1}{2} \right)^{n-2}\).
\end{solution}
\begin{problem}
如图，已知点 \(P\) 是 \(y\) 轴左侧（不含 \(y\) 轴）一点，抛物线 \(C:y^2=4x\) 上存在不同的两点 \(A,B\) 满足 \(PA,PB\) 的中点均在 \(C\) 上.

\begin{enumerate}
  \item 设 \(AB\) 中点为 \(M\)，证明：\(PM\) 垂直于 \(y\) 轴；
  \item 若 \(P\) 是半椭圆 \(x^2+\dfrac{y^2}{4}=1\)（\(x<0\)） 上的动点，求 \(\triangle PAB\) 面积的取值范围.
\end{enumerate}

\bitmapfigure[max width=0.42\linewidth,max totalheight=4.8cm]{img/2018/zhejiang/q21_fig1.png}

\end{problem}

\begin{solution}
（1）设 \(P(x_0,y_0)\)，\(A\left( \dfrac{1}{4}y_1^2,y_1 \right)\)，\(B\left( \dfrac{1}{4}y_2^2,y_2 \right)\).

因为 \(PA\)，\(PB\) 的中点在抛物线上，所以 \(y_1\)，\(y_2\) 为方程 \(\left( \frac{y+y_0}{2} \right)^2=4\cdot \frac{\dfrac{1}{4}y^2+x_0}{2}\)，即 \(y^2-2y_0y+8x_0-y_0^2=0\)
的两个不同的实数根.

所以 \(y_1+y_2=2y_0\).

因此，\(PM\) 垂直于 \(y\) 轴.

（2）由（1）可知
 \(\begin{cases}
y_1+y_2=2y_0,\\
y_1y_2=8x_0-y_0^2,
\end{cases}\)
所以
 \(|PM|=\frac{1}{8}\left( y_1^2+y_2^2 \right)-x_0=\frac{3}{4}y_0^2-3x_0\)，\(|y_1-y_2|=2\sqrt{2\left( y_0^2-4x_0 \right)}\).

因此，
 \(S_{\triangle PAB}=\frac{1}{2}|PM|\cdot |y_1-y_2|=\frac{3\sqrt{2}}{4}\left( y_0^2-4x_0 \right)^{\frac{3}{2}}\).

因为 \(x_0^2+\dfrac{y_0^2}{4}=1\)（\(x_0<0\)），所以
 \(y_0^2-4x_0=-4x_0^2-4x_0+4\in[4,5]\).

因此，\(\triangle PAB\) 面积的取值范围是
 \(\left[ 6\sqrt{2},\frac{15\sqrt{10}}{4} \right]\).
\end{solution}
\begin{problem}
已知函数 \(f(x)=\sqrt{x}-\ln x\).

\begin{enumerate}
  \item 若 \(f(x)\) 在 \(x=x_1,x_2\)（\(x_1\ne x_2\)） 处导数相等，证明：\(f(x_1)+f(x_2)>8-8\ln2\)；
  \item 若 \(a\le 3-4\ln2\)，证明：对于任意 \(k>0\)，直线 \(y=kx+a\) 与曲线 \(y=f(x)\) 有唯一公共点.
\end{enumerate}

\end{problem}

\begin{solution}
（1）函数 \(f(x)\) 的导函数
 \(f'(x)=\frac{1}{2\sqrt{x}}-\frac{1}{x}\).
由 \(f'(x_1)=f'(x_2)\) 得
 \(\frac{1}{2\sqrt{x_1}}-\frac{1}{x_1}=\frac{1}{2\sqrt{x_2}}-\frac{1}{x_2}\).

因为 \(x_1\ne x_2\)，所以
 \(\frac{1}{\sqrt{x_1}}+\frac{1}{\sqrt{x_2}}=\frac{1}{2}\).

由基本不等式得
 \(\frac{1}{2}\sqrt{x_1x_2}=\sqrt{x_1}+\sqrt{x_2}\ge2\sqrt[4]{x_1x_2}\).

因为 \(x_1\ne x_2\)，所以
 \(x_1x_2>256\).

由题意得
 \(f(x_1)+f(x_2)=\sqrt{x_1}-\ln x_1+\sqrt{x_2}-\ln x_2=\frac{1}{2}\sqrt{x_1x_2}-\ln(x_1x_2)\).

设
 \(g(x)=\frac{1}{2}\sqrt{x}-\ln x\)，
则
 \(g'(x)=\frac{1}{4x}(\sqrt{x}-4)\).
所以 \(g(x)\) 在 \([256,+\infty)\) 上单调递增，

故
 \(g(x_1x_2)>g(256)=8-8\ln2\)，
即
 \(f(x_1)+f(x_2)>8-8\ln2\).

（2）令
 \(m=\e^{-\left( |a|+k \right)}\)，\(n=\left( \frac{|a|+1}{k} \right)^2+1\)，
则
 \(f(m)-km-a>|a|+k-k-a\ge0\)，
 \(f(n)-kn-a<n\left( \frac{1}{\sqrt{n}}-\frac{a}{n}-k \right)\le n\left( \frac{|a|+1}{\sqrt{n}}-k \right)<0\)，
所以，存在 \(x_0\in(m,n)\) 使 \(f(x_0)=kx_0+a\).

所以，对于任意的 \(a\in \R\) 及 \(k\in(0,+\infty)\)，直线 \(y=kx+a\) 与曲线 \(y=f(x)\) 有公共点.

由 \(f(x)=kx+a\) 得
 \(k=\frac{\sqrt{x}-\ln x-a}{x}\).

设
 \(h(x)=\frac{\sqrt{x}-\ln x-a}{x}\)，
则
 \(h'(x)=\frac{\ln x-\dfrac{\sqrt{x}}{2}-1+a}{x^2}=\frac{-g(x)-1+a}{x^2}\)，
其中
 \(g(x)=\frac{\sqrt{x}}{2}-\ln x\).

由（1）可知 \(g(x)\ge g(16)\)，又 \(a\le3-4\ln2\)，

故
 \( -g(x)-1+a\le-g(16)-1+a=-3+4\ln2+a\le0\).
所以 \(h'(x)\le0\)，即函数 \(h(x)\) 在 \((0,+\infty)\) 上单调递减，因此方程 \(f(x)-kx-a=0\) 至多 \(1\) 个实根.

综上，当 \(a\le3-4\ln2\) 时，对于任意 \(k>0\)，直线 \(y=kx+a\) 与曲线 \(y=f(x)\) 有唯一公共点.
\end{solution}
